%! AP Statistics — Unit 9, Lesson 9.6 — Warm-Up (Answer Key)
\documentclass[10pt]{article}
\usepackage{apstats-article}
\usepackage{apstats-key}

\begin{document}

\pageheader{Unit 9, Lesson 9.6}{Warm-Up --- Answer Key}
\namedateperiod

\vspace{0.1in}

\textbf{Directions:} For each scenario below, answer (a), (b), and (c).

\begin{notesbox}{Scenario 1}
A health researcher randomly selects 200 adults and records whether each person
exercises regularly (yes or no). She wants to estimate the true proportion of adults
who exercise regularly.

\vspace{0.1in}
\begin{enumerate}[label=(\alph*), itemsep=6pt]
  \item Is the variable categorical or quantitative? \ans{categorical}
  \item How many groups or variables are involved? \ans{1 group / 1 proportion}
  \item Name the appropriate inference procedure: \ans{1-proportion $z$-interval}
\end{enumerate}
\end{notesbox}

\begin{notesbox}{Scenario 2}
A teacher randomly assigns 30 students to two groups. Group 1 uses a new study app;
Group 2 uses traditional methods. She wants to test whether the two groups differ in
mean exam score.

\vspace{0.1in}
\begin{enumerate}[label=(\alph*), itemsep=6pt]
  \item Is the variable categorical or quantitative? \ans{quantitative}
  \item How many groups are involved, and are they independent or paired? \ans{2 independent groups}
  \item Name the appropriate inference procedure: \ans{2-sample $t$-test for $\mu_1 - \mu_2$}
\end{enumerate}
\end{notesbox}

\begin{notesbox}{Scenario 3}
A biologist records the body length (cm) and weight (g) of 45 randomly selected fish
from a lake. She wants to test whether body length is a useful linear predictor of
weight.

\vspace{0.1in}
\begin{enumerate}[label=(\alph*), itemsep=6pt]
  \item Is the variable categorical or quantitative? \ans{quantitative}
  \item Describe the data structure (one mean, two means, bivariate, etc.): \ans{bivariate quantitative (slope)}
  \item Name the appropriate inference procedure: \ans{$t$-test for slope}
\end{enumerate}
\end{notesbox}

\end{document}
